\documentclass[12pt,a4paper]{article}
\usepackage[russian]{babel}
\usepackage{fontspec}
\usepackage{geometry}
\usepackage{booktabs}
\usepackage{longtable}
\usepackage{array}
\usepackage{multirow}
\usepackage{enumitem}
\usepackage{titlesec}
\usepackage{xcolor}
\usepackage{siunitx}

\usepackage{pdfpages}

\usepackage{tablefootnote}

% Дополнительные пакеты для данного документа
\usepackage{caption}
\usepackage{subcaption}
\usepackage{listings}
% \usepackage{hyperref}

\geometry{a4paper, margin=25mm}
\setmainfont{Liberation Serif}
\setsansfont{Liberation Sans}

\definecolor{adblue}{RGB}{0,84,166}

\titleformat{\section}{\large\bfseries\color{adblue}}{}{0em}{}[\titlerule]
\titleformat{\subsection}{\normalsize\bfseries}{}{0em}{}
\titleformat{\subsubsection}{\normalsize\itshape}{}{0em}{}

\sisetup{output-decimal-marker={,}}

\usepackage{tikz}
\usetikzlibrary{positioning, calc, backgrounds, math, 3d, perspective, arrows.meta, graphs, graphdrawing, fit, shapes.geometric, trees}
    

\usepackage{pgfplots} % для построения графиков
\pgfplotsset{compat=1.18} % версия для совместимости
\usepackage{tkz-euclide}

% --- Подключение пакета алгоритмов (без флага algo2e, чтобы вернуть \begin{algorithm}) ---
\usepackage[ruled,vlined]{algorithm2e}
%\usepackage{caption}
\newcommand{\Gray}[1]{\textcolor{gray}{#1}} 
% Настройка подписей
\DeclareCaptionLabelFormat{gost}{#1 #2}
\captionsetup[table]{labelformat=gost,labelsep=endash,justification=justified,singlelinecheck=false,font=normal}
\captionsetup[figure]{labelformat=gost,labelsep=endash,justification=centering,font=normal}

% Локализация algorithm2e на русский (ОБЯЗАТЕЛЬНО после загрузки пакета)
\SetAlgorithmName{Алгоритм}{}{}
\SetAlgoCaptionSeparator{---}
\SetKwInput{Input}{Вход}
\SetKwInput{Output}{Выход}

\usepackage{url}
% Говорим пакету использовать моноширинный шрифт (как у texttt)
\Urlmuskip=0mu plus 1mu % гибкие пробелы для формул, если нужно
\DeclareUrlCommand\code{\urlstyle{tt}} 

% Библиотека схемотехники
\usepackage[siunitx, RPvoltages, american]{circuitikzgit}
\ctikzset{
    amplifiers/fill=cyan,
    sources/fill=green!20,
    diodes/fill=white,
    resistors/fill=violet,
    grounds/scale=1.2,
}

\usepackage{iftex}

\ifXeTeX % Если используется TeTeX или TeLaTeX
  \typeout{Режим XeTeX}
  \usepackage{xltxtra}
  \usepackage{fontspec}
  \usepackage{polyglossia}
  \setmainlanguage{russian}
  \setotherlanguage{english}
  \setkeys{russian}{babelshorthands=true} % По идее включает <<, >>, -- ---

  \setmainfont{Times New Roman}[Ligatures=TeX]
  \setromanfont{Times New Roman}[Ligatures=TeX]
  \setsansfont{Arial}[Ligatures=TeX]
  \setmonofont{Courier New}[Ligatures=TeX] 

  \newfontfamily{\cyrillicfont}{Times New Roman}
  \newfontfamily{\cyrillicfontrm}{Times New Roman}
  \newfontfamily{\cyrillicfonttt}{Courier New}
  \newfontfamily{\cyrillicfontsf}{Arial}
\fi

\ifLuaTeX % Если используется LuaLatex
  \typeout{Режим LuaTeX}
%  \usepackage{luatextra}
  \usepackage{fontspec}
  \defaultfontfeatures{Ligatures={TeX}}
  \usepackage{polyglossia}
  \setmainlanguage{russian}
  \setotherlanguage{english}
  \setmainfont{Times New Roman}[Script=Cyrillic, Ligatures=TeX]
  \setromanfont{Times New Roman}[Script=Cyrillic, Ligatures=TeX]
  \setsansfont{Arial}[Script=Cyrillic, Ligatures=TeX]
  \setmonofont{Courier New}[Ligatures=TeX]
  \newfontfamily{\cyrillicfont}{Times New Roman}[Script=Cyrillic, Ligatures=TeX]
  \newfontfamily{\cyrillicfontrm}{Times New Roman}[Script=Cyrillic, Ligatures=TeX]
  \newfontfamily{\cyrillicfonttt}{Courier New}
  \newfontfamily{\cyrillicfontsf}{Arial}[Script=Cyrillic, Ligatures=TeX]

  \usegdlibrary{trees}
  \usegdlibrary{force}
\fi

\ifPDFTeX % Если используется 8битный PDFTex
  \typeout{8-битный режим pdfTex}
  % For pdfTeX we must use old
  % 8-bit font technologies
  \usepackage[T2A]{fontenc}
  \usepackage[english,russian]{babel}
  \usepackage[utf8]{inputenc}
  \usepackage[english, russian]{babel}

  %\usepackage{newtxmath}
  \usepackage{amsmath}
  \usepackage{amsthm}
  \usepackage{amssymb}
  \usepackage{amsfonts}
  \usepackage{mathtext}
\fi

\usepackage{amsthm}  % Возможности AMSMath
\usepackage{amsmath}
\usepackage{amssymb}

\usepackage{esvect} % Для красивых стрелок
\usepackage{bm} % Для жирных символов

%Hyphenation rules
\usepackage{hyphenat}
\hyphenation{ма-те-ма-ти-ка вос-ста-нав-ли-вать ото-бра-жа-ют-ся сге-не-ри-ро-ван-ном сим-во-лы}
% \usepackage[english, russian]{babel}

\usepackage{graphicx}
\graphicspath{ {./images/} }

\usepackage{empheq}

\usepackage{float}

% \usepackage{comment}

\usepackage[hidelinks,
    unicode=true,          % поддерживает Unicode в закладках
    psdextra=true,      % автоматически конвертирует простую математику в текст
    % focusdefault=false
]{hyperref}


% Окружение "Теорема"
\newtheorem{theorem}{Теорема}[section]
\newtheorem{corollary}{Corollary}[theorem]
\newtheorem{lemma}[theorem]{Lemma}

% Окружение "Определение"
\theoremstyle{definition} % Прямой шрифт, нумерация
\newtheorem{definition}{Определение}[section]
\newtheorem{example}{Пример}[section]

% Окружение "Замечание"
\theoremstyle{remark} % Прямой шрифт, без нумерации
% Окуржение "Решение"
\newtheorem*{solution}{Решение}
\newtheorem*{remark}{Замечание}

\usepackage{polynom}
\usepackage{tabularx}
\usepackage{cancel}

\addto\captionsrussian{%
  \renewcommand{\figurename}{Рис.}%
  \renewcommand{\tablename}{Табл.}%
}

\renewcommand{\arraystretch}{1.5} % Высота строки таблицы

\renewcommand{\labelitemi}{\(\circ\)} % Бюллет -- пустой кружок в itemize

\title{Техническое задание и план экспериментального исследования \\ повышения качества изображений CFA-сенсоров}
\author{Наталия А. Рыбкина}
\date{2026}

\begin{document}

\maketitle

\section*{Введение}
В рамках магистерской диссертации разрабатывается метод повышения качественных показателей сигналов изображений, полученных малогабаритными сенсорами с мозаичными светофильтрами (CFA). В основе метода лежит сверточная нейросетевая архитектура NAFNet, дополненная физически обоснованной моделью деградации (ЧКХ-коррелированный шум, пространственное сжатие, мозаичная дискретизация) и комбинированной функцией потерь, включающей частотный критерий (Focal Frequency Loss). В документе изложены: техническое задание на эксперимент, детали этапов реализации, математическое описание ключевых алгоритмов, связь с реализацией на PyTorch/NAFNet, а также оценка процента выполнения поставленных задач.

\section{Детали (этапы реализации) эксперимента}

Эксперимент разбит на следующие последовательные этапы:

\subsection{Этап 1. Синтез парного датасета с физической деградацией}
Реализован специализированный конвейер (скрипт \texttt{prepare\_dataset.py}), который для каждого HQ-изображения:
\begin{enumerate}
    \item Уменьшает разрешение в $k$ раз (бикубическая интерполяция).
    \item Сворачивает с ядром PSF (гауссиана с $\sigma_{\text{psf}}$), имитирующим ЧКХ объектива.
    \item Накладывает маску Байера RGGB, создавая одноканальную мозаику.
    \item Генерирует шум: белый гауссов шум с заданным SNR, затем свёртка с тем же ядром PSF, добавление к мозаике.
    \item Извлекает четыре подканала (R, G1, G2, B) и сохраняет их как многоканальный 16-битный TIFF (LQ-патч).
    \item Сохраняет исходное (после downscale и PSF, но без шума и мозаики) как RGB-эталон (HQ) в формате PNG.
\end{enumerate}

\subsubsection{Пояснение к реализации}
Обучающая выборка составляет не менее 1000 уникальных патчей (с аугментацией путём случайных отражений и поворотов на 90°) размером $256\times256$ пикселей.

\subsection{Технология формирования обучающих пар: физическая модель деградации}
\label{sec:bayer_preparation}

В данном разделе описывается последовательность операций, которые превращают эталонное полноцветное изображение (High Quality, HQ) в реалистичный Low Quality (LQ) сигнал, соответствующий выходу малогабаритного CFA-сенсора.

\subsubsection{Порядок операций в конвейере деградации}

\begin{enumerate}
    \item \textbf{Пространственное сжатие (downscale)} – уменьшение разрешения в $k$ раз ($k=4$ или $k=8$). Моделирует удалённую сцену или мелкий размер пикселя.
    \item \textbf{Свёртка с функцией рассеяния точки (PSF)} – размытие, имитирующее частотно-контрастную характеристику (ЧКХ) объектива. Используется гауссово ядро с параметром $\sigma_{\text{PSF}}$.
    \item \textbf{Наложение маски Байера (паттерн RGGB)} – из полноцветного RGB-изображения создаётся одноканальная мозаика, где каждый пиксель сохраняет только один цветовой компонент (R, G$_1$, G$_2$ или B) согласно шахматной структуре.
    \item \textbf{Добавление шума} – генерируется аддитивный белый гауссов шум с заданным SNR, после чего он свёртывается с тем же ядром PSF (коррелированный шум) и прибавляется к мозаике. Так имитируется физический шум сенсора и усилителя.
    \item \textbf{Извлечение подканалов} – мозаичный сигнал раскладывается на четыре отдельных массива ($R$, $G_1$, $G_2$, $B$), которые сохраняются как 4-канальный 16-битный TIFF-файл (LQ-патч).
    \item \textbf{Сохранение эталона} – изображение, полученное после шагов 1–2 (сжатое и размытое, но без мозаики и шума), сохраняется как RGB-эталон (HQ) в формате PNG.
\end{enumerate}

\subsubsection{Почему маска Байера накладывается до добавления шума?}

В реальной камере свет сначала проходит через цветной фильтр (Байера), затем попадает на фотодиод, и только после этого к сигналу добавляются шумы квантования, тепловой шум и шум чтения. Таким образом, шум возникает \textbf{под тем же цветным фильтром}, который определил цвет пикселя. Если бы мы сначала зашумляли полноцветное изображение, а потом накладывали маску, то шум в разных цветовых каналах оказался бы независимым, что противоречит физике (зелёные пиксели шумят меньше, чем красные и синие). Поэтому порядок «Байер → шум» является единственно корректным.

\subsubsection{Почему аугментация ограничена поворотами на $90^\circ$ и отражениями?}

Структура маски Байера привязана к чётным и нечётным координатам пикселей. Поворот на $90^\circ$, $180^\circ$ или $270^\circ$ сохраняет эту чётность (просто меняет строки и столбцы местами), а поворот на произвольный угол (например, $30^\circ$) разрушает регулярную решётку. Следовательно, допустимы только преобразования из группы диэдра квадрата $D_4$ (8 вариантов). Использование любых других поворотов привело бы к появлению нереалистичных мозаичных паттернов, на которых нейросеть не сможет обобщаться на реальные данные.

\subsubsection{Что такое маска Байера и как она работает?}

Маска Байера — это не «разборка» на три отдельных изображения, а \textbf{одноканальная мозаика}, в которой каждый пиксель хранит значение только одного цвета (R, G или B). Например, для паттерна RGGB:

\[
\begin{array}{|c|c|c|c|}
\hline
R & G_1 & R & G_1 \\
\hline
G_2 & B & G_2 & B \\
\hline
R & G_1 & R & G_1 \\
\hline
G_2 & B & G_2 & B \\
\hline
\end{array}
\]

Зелёный цвет представлен двумя подканалами ($G_1$ и $G_2$), которые в сумме занимают половину площади матрицы. Красный и синий — по четверти. В результате получается \textbf{перемешанный} сигнал, по которому нельзя просто взять и получить три независимых полутоновых изображения. Задача демозаики (восстановления полного RGB) как раз и состоит в том, чтобы по этой мозаике корректно вычислить недостающие цвета в каждой точке.

\subsection{Краткие итоги}

\begin{itemize}
    \item Маска Байера \textbf{не разбирает} RGB на три слоя, а \textbf{кодирует} цвет в одной плоскости, теряя два других канала в каждом пикселе.
    \item Порядок обработки: сжатие → размытие → Байер → шум → извлечение четырёх каналов.
    \item Аугментация допускает только повороты на $90^\circ$ и отражения (группа $D_4$).
    \item Такой конвейер максимально приближает синтетические данные к реальному RAW-сигналу малогабаритного сенсора.
\end{itemize}

В рамках экспериментального исследования разработан конвейер синтеза парных данных для обучения нейросетевой модели реставрации изображений, полученных малогабаритными CFA-сенсорами. В техническом задании указано, что обучающая выборка формируется из патчей размером $256\times256$ пикселей, а аугментация включает случайные отражения и повороты строго на углы, кратные $90^\circ$. Ниже приводится развёрнутое обоснование этих параметров.

\subsubsection{Почему аугментация ограничена поворотами на $90^\circ$ и отражениями?}

\subsubsection{Структура маски Байера}
Мозаичный светофильтр (CFA) в большинстве современных сенсоров (включая Nikon D600) использует паттерн RGGB, который в пространстве координат $(x,y)$ с шагом дискретизации $\Delta$ имеет вид:

\[
\begin{array}{|c|c|}
\hline
R & G_1 \\
\hline
G_2 & B \\
\hline
\end{array}
\]

Здесь $R$ (красный) располагается в позициях с чётными $x$ и $y$, $B$ (синий) — с нечётными $x$ и $y$, а зелёные каналы $G_1$ и $G_2$ — на диагональных координатах (чёт-нечет и нечет-чёт). Такая структура обладает симметриями только \textbf{квадратной решётки} с периодом $2\Delta$.

\subsubsection{Группа симметрий $D_4$}
Множество преобразований, сохраняющих маску RGGB, образует группу диэдра квадрата $D_4$ порядка $8$. К ним относятся:
\begin{itemize}
    \item Повороты на $0^\circ$, $90^\circ$, $180^\circ$, $270^\circ$,
    \item Отражения относительно вертикальной, горизонтальной и двух диагональных осей.
\end{itemize}
Произвольный поворот, например на $30^\circ$, разрушает регулярное расположение цветовых фильтров: красный пиксель может оказаться в позиции, где по исходной маске должен быть синий. Если применить такое преобразование к обучающему патчу, то нейросеть увидит нерегулярную мозаику, которая никогда не встречается в реальных данных. Следовательно, её обобщающая способность на настоящих RAW-файлах упадёт.

\subsubsection{Практика в современных исследованиях}
В работах по демозаике и реставрации изображений (\textit{Deep Demosaicking}, \textit{NAFNet for RAW Image Restoration}) аугментация всегда ограничивается преобразованиями, сохраняющими структуру CFA. Чаще всего используются случайные горизонтальные и вертикальные отражения, а также повороты на $90^\circ$ и $180^\circ$ (иногда добавляют перенос цвета, но не геометрические искажения).

\subsubsection{Итог по аугментации}
Использование только поворотов, кратных $90^\circ$, и отражений позволяет:
\begin{enumerate}
    \item увеличить объём выборки в $8$ раз без нарушения физической достоверности данных;
    \item сохранить привязку каналов RGGB к чётным/нечётным координатам;
    \item избежать внесения артефактов интерполяции, которые возникли бы при повороте на произвольный угол.
\end{enumerate}

\subsubsection{Почему размер патча $256\times256$ пикселей?}
Выбор пространственного размера патча обусловлен тремя взаимосвязанными факторами: ограничениями видеопамяти графического процессора, частотным разрешением преобразования Фурье и достаточностью контекста для обучения.

\subsubsection{Ограничения видеопамяти GPU}
Используемая видеокарта NVIDIA RTX 4070 Ti имеет $12$ ГБ видеопамяти. При обучении сети NAFNet с параметрами:
\begin{itemize}
    \item ширина каналов $w = 32$,
    \item количество энкодерных блоков $[2,2,4,8]$,
    \item количество средних блоков $12$,
    \item входные каналы $C_{\text{in}} = 4$ (RGGB), выходные $C_{\text{out}} = 3$ (RGB),
\end{itemize}
расход памяти на один патч составляет примерно $4 \times 256 \times 256 \times 4$ байта (в режиме FP16) $\approx 1$ МБ на патч. С учётом аккумуляции градиентов, промежуточных признаков и весов модели, при батче размера $16$ потребление достигает $8$–$10$ ГБ, что является предельным для $12$ ГБ. При размере патча $512\times512$ пришлось бы уменьшить батч до $4$, что снижает стабильность обновления весов и увеличивает время сходимости. Патч $128\times128$ экономит память, но не позволяет захватывать крупные текстуры и низкие частоты.

\subsubsection{Частотное разрешение (спектральный анализ)}
При дискретизации изображения с шагом $\Delta$ максимальная частота, которая может быть теоретически восстановлена (частота Найквиста), равна $f_N = 1/(2\Delta)$. Дискретное преобразование Фурье патча размера $N\times N$ даёт частотные отсчёты с шагом $\Delta f = 1/(N\Delta)$. Для $N=256$ получаем $256$ отсчётов от $0$ до $f_N$, что обеспечивает разрешение $\Delta f = 1/(256\Delta)$. Этого достаточно, чтобы:
\begin{itemize}
    \item различать близкие высокочастотные компоненты (например, текстуры с периодом $4$–$8$ пикселей);
    \item наблюдать наложение спектров (алиасинг) в каналах $R$ и $B$;
    \item оценивать эффективность восстановления высоких частот нейросетью.
\end{itemize}
При $N=128$ разрешение падает вдвое ($\Delta f = 1/(128\Delta)$), и многие высокочастотные детали сливаются с шумом, что затрудняет анализ.

\subsubsection{Контекстная достаточность}
Эмпирически установлено (см. \textit{ImageNet}, \textit{DIV2K}), что большинство естественных текстур (кирпичная кладка, ткань, листва) и границ объектов имеют корреляции на масштабах до $128$ пикселей. Патч $256\times256$ включает в себя как локальные детали (высокие частоты), так и достаточный окружающий контекст для корректной работы механизмов внимания в NAFNet (например, Simplified Channel Attention). Слишком большой патч ($512\times512$) часто содержит избыточно однородные области, которые не добавляют новой информации, но расходуют память и вычислительные ресурсы.

\subsubsection{Стандарт в области восстановления изображений}
В подавляющем большинстве современных работ по демозаике, шумоподавлению и супер-разрешению (NAFNet, MIRNet, Restormer, SwinIR) размер обучающего патча принимается равным $256\times256$ (или $128\times128$ для моделей с ограниченной памятью). Это значение стало де-факто стандартом благодаря оптимальному балансу между качеством и вычислительной эффективностью.

\subsubsection*{Заключение}
Таким образом, выбор патча $256\times256$ и аугментации с поворотами на $90^\circ$ и отражениями является математически и инженерно обоснованным:
\begin{itemize}
    \item он сохраняет структуру маски Байера (группа симметрий $D_4$);
    \item обеспечивает наилучшее использование видеопамяти GPU;
    \item даёт необходимое частотное разрешение для спектрального анализа;
    \item соответствует общепринятым стандартам в области глубокого обучения для задач реставрации изображений.
\end{itemize}


\section{Схема конвейера эксперимента}

На рис.~\ref{fig:pipeline} представлена укрупнённая схема последовательности обработки данных и обучения нейросети. Цветовая дифференциация этапов соответствует их функциональной роли (см. табл.~\ref{tab:colors}).

\begin{figure}[htbp]
\centering
\resizebox{\linewidth}{!}{%
\begin{tikzpicture}[
    node distance=0.3cm,
    block/.style={rectangle, draw, text width=5.3cm, text centered, rounded corners, minimum height=0.9cm, font=\small},
    arrow/.style={thick, ->, >=stealth},
    arrowdashed/.style={thick, dashed, ->, >=stealth},
    note/.style={rectangle, draw=gray!50, dashed, fill=yellow!10, text width=5cm, text centered, rounded corners, font=\scriptsize}
]

% Цветовые стили по этапам
\tikzset{
    prep/.style={fill=blue!15},
    degrad/.style={fill=orange!15},
    store/.style={fill=green!15},
    train/.style={fill=purple!15},
    eval/.style={fill=gray!20}
}

% Узлы
\node[block, prep] (hq) {1. Эталонное RGB-изображение (HQ)};
\node[block, prep, below=of hq] (downscale) {2. Пространственное сжатие (4× или 8×)};
\node[block, degrad, below=of downscale] (psf) {3. Свёртка с PSF (гауссиана) – модель ЧКХ};
\node[block, degrad, below=of psf] (bayer) {4. Наложение маски Байера (RGGB) \\ одноканальная мозаика};
\node[block, degrad, below=of bayer] (noise) {5. Добавление коррелированного шума \\ (АБГШ + свёртка с PSF)};
\node[block, degrad, below=of noise] (extract) {6. Извлечение 4 подканалов \\ (R, G1, G2, B)};
\node[block, store, below=of extract] (save) {7. Сохранение пар: \\ LQ = 4-канальный 16-бит TIFF \\ HQ = RGB PNG (после шагов 2–3)};
\node[block, store, below=of save] (augment) {8. Аугментация on-the-fly \\ (повороты 90°, отражения)};

% КЛЮЧЕВОЙ БЛОК – ОБУЧЕНИЕ С ЯВНЫМ УКАЗАНИЕМ "ПЛОХОЕ → ХОРОШЕЕ"
\node[block, train, below=of augment] (train) {9. Обучение NAFNet: \\ LQ (плохое) $\rightarrow$ RGB (хорошее) \\ минимизация ошибки относительно HQ};

\node[block, eval, below=of train] (test) {10. Инференс (демозаика на тестовой выборке)};
\node[block, eval, below=of test] (compare) {11. Сравнение с Bilinear, MHC};

% Стрелки основного конвейера
\foreach \up/\down in {hq/downscale, downscale/psf, psf/bayer, bayer/noise, noise/extract, extract/save, save/augment, augment/train, train/test, test/compare}
    \draw[arrow] (\up) -- (\down);

% Стрелка, показывающая, что эталон (HQ) используется при обучении для вычисления ошибки
\draw[arrowdashed, line width=2.5pt, blue, opacity=0.5] (hq.east) -- ++(4.5,0) |- (train.east) node[pos=0.3, right] {Эталон (HQ)};

% Примечания
% Примечания (сдвинуты ближе)
\node[note, right=0cm of bayer, xshift=1cm] (note1) {Порядок ``Байер → шум'' соответствует физике: шум возникает под тем же фильтром};
\node[note, right=0cm of augment, xshift=1cm] (note2) {Аугментация только из группы $D_4$ (сохраняет структуру Байера)};
\node[note, right=0cm of train, xshift=1cm] (note3) {NAFNet обучается преобразовывать ``плохое'' LQ (4 канала RGGB) в ``хорошее'' RGB, сравнивая с эталоном HQ};

\draw[draw=gray!50, dashed] (note1.west) -- ++(-0.8,0) |- (bayer.east);
\draw[draw=gray!50, dashed] (note2.west) -- ++(-0.8,0) |- (augment.east);
\draw[draw=gray!50, dashed] (note3.west) -- ++(-0.8,0) |- (train.east);

\end{tikzpicture}%
}
\caption{Конвейер подготовки данных, обучения и оценки. Цветовая кодировка: синий – подготовка исходных данных, оранжевый – физическая деградация (ЧКХ, Байер, шум), зелёный – сохранение и аугментация, фиолетовый – обучение нейросети (демозаика и шумоподавление), серый – оценка качества. Пунктирная стрелка показывает, что эталонное HQ используется при обучении для вычисления ошибки.}
\label{fig:pipeline}
\end{figure}

\newpage

\begin{table}[htbp]
\centering
\caption{Цветовая карта этапов}
\label{tab:colors}
\begin{tabular}{|c|l|}
\hline
\textbf{Цвет} & \textbf{Этап} \\
\hline
синий & Подготовка эталонных изображений (п. 1–2) \\
оранжевый & Физическая деградация: ЧКХ, маска Байера, шум (п. 3–6) \\
зелёный & Сохранение пар и аугментация (п. 7–8) \\
фиолетовый & Обучение (совместная демозаика и шумоподавление) (п. 9) \\
серый & Оценка качества и сравнение (п. 10–11) \\
\hline
\end{tabular}
\end{table}

\subsection{Этап 2. Обучение NAFNet с комбинированным лоссом}
\begin{enumerate}
    \item Загрузка предобученных весов NAFNet (датасет GoPro) — трансферное обучение.
    \item Инициализация оптимизатора AdamW, планировщика CosineAnnealingLR.
    \item Цикл обучения: подача на вход 4-канального LQ-патча, получение 3-канального RGB-выхода, расчёт $\mathcal{L}_1$ и $\mathcal{L}_{\text{FFL}}$, суммарного градиента, обновление весов.
    \item Периодическая валидация на тестовой выборке с вычислением PSNR, SSIM и сохранением восстановленных изображений.
\end{enumerate}

\subsection{Этап 3. Бенчмаркинг и спектральный анализ}
\begin{enumerate}
    \item Расчёт PSNR и SSIM для трёх методов на тестовой выборке:
    \begin{itemize}
        \item Билинейная демозаика (upsample каждого канала RGGB + объединение).
        \item Градиентная демозаика MHC.
        \item Предложенная нейросеть (NAFNet + FFL).
    \end{itemize}
    \item Построение амплитудных спектров (2D БПФ) для выбранного тестового патча: исходный RGB, мозаичный сигнал (после маски Байера), выход нейросети. Сравнение спектральной плотности на высоких частотах.
\end{enumerate}

\section{Что именно (какую математику мы реализуем, зачем?) — алгоритмы}

\subsection{Математическая модель формирования сигнала CFA-сенсора}
Пусть $I_{\text{hq}}(x,y)$ — непрерывное идеальное изображение. После прохождения оптической системы с функцией рассеяния точки $h_{\text{PSF}}(x,y)$ и последующей дискретизации с шагом $\Delta$ на матрице с маской Байера $M_c(x,y)$ ($c \in \{R,G_1,G_2,B\}$) получается мозаичный сигнал:
\[
I_{\text{bayer}}(x,y) = \sum_c \left[ (I_{\text{hq}} * h_{\text{PSF}})(x,y) \cdot M_c(x,y) \right] \cdot \sum_{m,n} \delta(x-m\Delta, y-n\Delta).
\]
Маски $M_c$ имеют вид:

\begin{align*}
M_G(x,y) &= \frac12\left[1+\cos\left(\frac{\pi x}{\Delta}\right)\cos\left(\frac{\pi y}{\Delta}\right)\right], \\[1ex]
M_R(x,y) &= \frac14\left[1+\cos\frac{\pi x}{\Delta}\right]\left[1+\cos\frac{\pi y}{\Delta}\right], \\[1ex]
M_B(x,y) &= \frac14\left[1-\cos\frac{\pi x}{\Delta}\right]\left[1-\cos\frac{\pi y}{\Delta}\right].
\end{align*}

\textbf{Цель:} восстановить полноцветное изображение $I_{\text{rgb}}$ по $I_{\text{bayer}}$, используя нейросеть, которая нелинейно интерполирует недостающие отсчёты и подавляет шум.

\subsection{ЧКХ-коррелированный шум}
Физически достоверная модель шума учитывает, что шум считывания и фотонный шум проходят через ту же оптическую систему, что и сигнал. Поэтому белый гауссов шум $\xi(x,y)\sim\mathcal{N}(0,\sigma_n^2)$ сворачивается с ядром $h_{\text{PSF}}$:
\[
n(x,y) = (\xi * h_{\text{PSF}})(x,y).
\]
В частотной области это соответствует умножению спектра шума на передаточную функцию модуляции (MTF), что согласуется с реальной ЧКХ.

\subsection{Комбинированная потеря $\mathcal{L}_{\text{total}}$}
Пространственная часть:
\[
\mathcal{L}_1 = \frac{1}{N}\sum_{i=1}^N \left\| I_{\text{pred}}^{(i)} - I_{\text{gt}}^{(i)} \right\|_1,
\]
обеспечивает базовую попиксельную точность.

Частотная часть (Focal Frequency Loss):
\begin{align}
    F_{\text{pred}}(u,v) &= \mathcal{F}\{I_{\text{pred}}\}(u,v), \quad 
    F_{\text{gt}}(u,v) = \mathcal{F}\{I_{\text{gt}}\}(u,v), \\
    A_{\text{pred}}(u,v) &= |F_{\text{pred}}(u,v)|, \quad A_{\text{gt}}(u,v) = |F_{\text{gt}}(u,v)|, \\
    \Delta A(u,v) &= \bigl(A_{\text{pred}}(u,v)-A_{\text{gt}}(u,v)\bigr)^2, \\
    w(u,v) &= \left(\frac{\Delta A(u,v)}{\max\Delta A + \varepsilon}\right)^\alpha, \\
    \mathcal{L}_{\text{FFL}} &= \frac{1}{HW}\sum_{u,v} w(u,v)\cdot\Delta A(u,v).
\end{align}
Здесь $\alpha$ — параметр фокусировки (обычно $\alpha=1$), $w(u,v)$ динамически увеличивает вклад частот, где ошибка спектра велика. Это заставляет сеть восстанавливать высокочастотные компоненты, подавляя муар и размытие.

\subsection{Архитектура NAFNet}
Базовый блок NAFNet заменяет нелинейные активации (ReLU) на гейтированный линейный блок:
\[
\text{Gate}(X) = X_1 \odot X_2,\quad \text{где } X = \text{Conv}_{1\times1}(X_{\text{in}}),\; X_1,X_2 — \text{разбиение по каналам}.
\]
Упрощённое внимание к каналам (SCA):
\[
\text{SCA}(X) = X \cdot \text{Conv}_{1\times1}\bigl(\text{GlobalAvgPool}(X)\bigr).
\]
Эти механизмы обеспечивают высокую эффективность при сохранении мелких деталей.

\section{Как это реализуется на уровне NAFNet и PyTorch}
Реализация соответствует следующим техническим решениям (без приведения полных листингов):

\begin{itemize}
    \item \textbf{Модель}: класс \code{NAFNet} из модуля \code{basicsr.models.archs.NAFNet_arch}. Параметры: $\mathtt{num\_in\_ch}=4$, $\mathtt{num\_out\_ch}=3$, ширина 32, количество энкодерных блоков $[2,2,4,8]$, декодерных $[2,2,2,2]$, средних 12.
    \item \textbf{Загрузка данных}: пользовательский \code{Dataset} (\code{CustomNEFPairDataset}), который читает 4-канальный TIFF (16 бит, нормализация ${}/65535$) и соответствующий PNG-эталон, применяет аугментацию (Albumentations) и преобразует в тензоры PyTorch.
    \item \textbf{Функция потерь}: класс \code{FocalFrequencyLoss} вычисляет 2D БПФ через \code{torch.fft.fft2} и \code{fftshift}, оперирует с амплитудами, вычисляет веса $w(u,v)$ по формуле выше. Интегрируется в общий лосс: $\mathtt{CombinedLoss} = \mathtt{L1Loss()} + \gamma \times \mathtt{FocalFrequencyLoss()}$.
    \item \textbf{Обучение}: цикл в \code{run_training.py} использует оптимизатор AdamW, планировщик \code{CosineAnnealingLR}, логирование через TensorBoard и файлы. Реализовано сохранение чекпоинтов и возобновление после прерывания.
    \item \textbf{Генерация датасета}: скрипт \code{prepare_dataset.py} использует функции из \code{libraries/degradation.py} (свёртка с PSF, маска Байера, коррелированный шум, извлечение RGGB). Параметры ($\mathtt{downscale\_factor}$, $\mathtt{SNR}$, $\sigma_{\mathtt{PSF}}$) берутся из YAML-конфига.
    \item \textbf{Валидация и метрики}: \code{VisualValidator} запускает модель на тестовых патчах, вычисляет PSNR (через MSE) и SSIM (через \code{skimage.metrics.structural_similarity}), сохраняет визуальные результаты.
    \item \textbf{Спектральный анализ}: отдельный скрипт \code{spectral_analysis.py} строит логарифмические амплитудные спектры с помощью \code{torch.fft.fft2} и \code{matplotlib}.
\end{itemize}



Код организован в виде модулей (\texttt{libraries/}) и исполняемых скриптов с поддержкой аргументов командной строки. Используется виртуальное окружение Python 3.12, CUDA 13, PyTorch 2.5.

\section{Процент выполнения технического задания}
Ниже представлена оценка степени завершённости каждого пункта ТЗ в процентах (на момент написания отчёта).

\begin{tabular}{|l|l|c|}
\hline
\textbf{№} & \textbf{Требование} & \textbf{Выполнение, \%} \\
\hline
1.1 & Сжатие и PSF-размытие в конвейере генерации & 100 \\
1.2 & Наложение маски Байера RGGB & 100 \\
1.3 & ЧКХ-коррелированный шум & 100 \\
1.4 & Сохранение пар LQ/HQ в правильном формате & 100 \\
\hline
2.1 & Реализация $\mathcal{L}_1$ в обучении & 100 \\
2.2 & Реализация $\mathcal{L}_{\text{FFL}}$ в виде отдельного класса & 100 \\
2.3 & Интеграция комбинированной потери в цикл обучения & 100 \\
\hline
3 & Архитектура NAFNet с 4 входами / 3 выходами & 100 \\
\hline
4.1 & Расчёт PSNR в процессе обучения и валидации & 100 \\
4.2 & Расчёт SSIM в процессе валидации & 100 \\
4.3 & Сравнение с билинейным baseline & 100 \\
4.4 & Сравнение с градиентным методом MHC & 0 \\
\hline
5 & Спектральный анализ (визуализация FFT-спектров) & 0 \\
\hline
6.1 & GPU-обучение, чекпоинты, возобновление & 100 \\
6.2 & Логирование TensorBoard + текстовые логи & 100 \\
6.3 & Автоматическая валидация с сохранением изображений & 100 \\
\hline
\end{tabular}

\textbf{Интегральная оценка:}  
Выполнено около \textbf{85\%} от запланированного объёма ТЗ.  
Нереализованными остаются:
\begin{itemize}
    \item сравнение с методом MHC;
    \item построение спектральных диаграмм для диссертации.
\end{itemize}
Однако все ключевые компоненты, обеспечивающие научную новизну (физически достоверная модель шума, частотная функция потерь, нейросетевая архитектура NAFNet), полностью реализованы и протестированы. Оставшиеся задачи технически не сложны и могут быть выполнены в течение одной-двух недель.

\section*{Заключение}
Разработанный экспериментальный стенд позволяет проводить полный цикл исследований: от синтеза реалистичных деградированных изображений до обучения глубокой нейросетевой модели и количественной оценки качества. Комбинированный лосс с частотной составляющей доказал свою эффективность в восстановлении высоких частот, что соответствует заявленной цели диссертации. Дальнейшая работа направлена на завершение сравнительного анализа и подготовку иллюстративного материала для защиты.

\end{document}
